\documentclass[11pt]{article}

\usepackage[margin=1in]{geometry}
\usepackage{amsmath,amssymb}
\usepackage{booktabs}
\usepackage{xcolor}
\usepackage{framed}
\usepackage{hyperref}
\usepackage{enumitem}

\definecolor{auditbg}{gray}{0.93}
\definecolor{auditbar}{gray}{0.4}
\definecolor{fixedbg}{RGB}{230,240,255}
\definecolor{fixedbar}{RGB}{80,110,180}

\newenvironment{auditbox}{%
  \def\FrameCommand{\fboxsep=8pt\fboxrule=0pt\colorbox{auditbg}{\vrule width 3pt\relax\color{auditbar}\vrule width 0pt}\colorbox{auditbg}}%
  \MakeFramed{\advance\hsize-2\width \FrameRestore}\noindent
  \textbf{Needs audit.}\ %
}{%
  \endMakeFramed\medskip
}

\newenvironment{fixedbox}{%
  \def\FrameCommand{\fboxsep=8pt\fboxrule=0pt\colorbox{fixedbg}{\vrule width 3pt\relax\color{fixedbar}\vrule width 0pt}\colorbox{fixedbg}}%
  \MakeFramed{\advance\hsize-2\width \FrameRestore}\noindent
  \textbf{Found and fixed.}\ %
}{%
  \endMakeFramed\medskip
}

\title{The training advisor: a specification of \texttt{advisor.py}}
\author{}
\date{Describes \texttt{src/training/advisor.py} against exercise dictionary version 9.\\
Spec: \texttt{docs/GUIDE.md} Part B. Parameter review: \texttt{docs/AUDIT.md}.}

\begin{document}
\maketitle
\tableofcontents

\begin{auditbox}
\textbf{This entire document describes a retired scoring model.} Before
any real training log existed, \texttt{advisor.py} was rebuilt around a
single continuous-decay EWMA per cell (\texttt{score = -log(ewma/k)}) and
one hard overdue clock, replacing the raw-deficit/window/staleness/
normalization/$\gamma$/resource-interference machinery this document
formalizes throughout. The resource-cost matrix $R$ and its interference
scoring are gone from the model entirely, not merely re-derived. See
\texttt{docs/GUIDE.md}'s Scoring and Overdue sections for the
live model, and \texttt{docs/STATUS.md} for why the reversal
happened before this document was updated to match. Every section below
this point --- the deficit derivation, the overdue-clock table, the
interference formula, the worked example --- describes dictionary version
7--9 and is kept only as a correct record of what that model did, not as
current documentation. A full rewrite against the live model is tracked
but not yet done.
\end{auditbox}

\section{What this is}

The advisor is a ranking function, not a program generator: it answers
``what is most needed today,'' and a human decides how many pieces to do
and designs the actual session. Nothing here is a claim about training
science beyond what \texttt{docs/GUIDE.md} states explicitly ---
this document only formalizes what the code computes.

Grey boxes throughout mark a value the code needs to run but that
\texttt{GUIDE.md} either left as a free choice or never specified at all
--- these are exactly the items in \texttt{docs/AUDIT.md}, placed
here at the point in the math where each one enters. Blue boxes mark a
bug that was found (via a two-month closed-loop simulation) and already
fixed; they document a correctness constraint the implementation must
satisfy, not an open question.

\section{Axes}

Two small finite sets index everything below.

\begin{align*}
\mathcal{P} &= \{\mathrm{hinge}, \mathrm{push}, \mathrm{squat}, \mathrm{pull}, \mathrm{flexion}, \mathrm{brace}, \mathrm{engine}\} & \text{(patterns, 7)}\\
\mathcal{S} &= \{\mathrm{light}, \mathrm{heavy}, \mathrm{hhr}\} & \text{(styles, 3)}
\end{align*}

\texttt{squat}, \texttt{hinge}, \texttt{push}, and \texttt{pull} name
themselves; \texttt{flexion} is loaded trunk/hip flexion through range
(archetype: toes-to-bar, V-up) and \texttt{brace} is resisting movement
under load (planks, loaded carries, the Turkish get-up). The last two
were a single \texttt{midline} pattern until dictionary version 8 ---
see \S\ref{sec:patternweight}. \texttt{engine} is restored in version 9 in
a deliberately narrow form: it contains \emph{only} exercises whose $E$
row would otherwise be empty (running, biking, rowing, double-unders,
burpees). That restriction is what makes it compatible with $E$'s
row-stochasticity --- since no member carries any other pattern weight,
crediting engine never forces a split that would dilute an existing
pattern. A general engine pattern would: giving a thruster engine credit
means subtracting it from $\mathrm{squat}$ and $\mathrm{push}$, because
rows sum to at most $1$ and weights split volume rather than duplicating
it. The first four are Pat Sherwood's CrossFit Linchpin gross movement
patterns (originally coded \texttt{BP}/\texttt{PFG}/\texttt{PR}/\texttt{UBP}/\texttt{ML})
renamed to match the ``squat/hinge/push/pull'' framing common in general
S\&C (Dan John and others). Sherwood's ``engine'' was originally dropped,
on the grounds that it classifies energy system rather than movement and
duplicated the modality axis --- the latter being the styles axis together
with the block's logged time-domain fields. Both objections stand against
a \emph{general} engine pattern, and neither survives the narrow one
above, which is why running and biking no longer sit in the residual.
\texttt{hhr} (``heavy
at high heart rate'') is a stimulus, not a session type: load carried
while breathing hard, roughly 60--70\% of 1RM for reps inside a workout.

\section{The exercise dictionary}

Let $\mathcal{E}$ be the set of exercises (63 entries in dictionary
version 9). Each exercise $e \in \mathcal{E}$ carries:

\begin{itemize}[nosep]
  \item $E_{e} : \mathcal{P} \to [0,1]$, the \textbf{pattern-weight row},
    with $\sum_{p \in \mathcal{P}} E_{e,p} \le 1$. The unallocated mass
    $1 - \sum_p E_{e,p}$ is residual --- volume the model doesn't
    attribute to any pattern. A squat clean is $0.5\,\mathrm{hinge} +
    0.5\,\mathrm{squat}$; a thruster is $0.5\,\mathrm{squat} + 0.5\,\mathrm{push}$;
    a clean \& jerk is $0.5\,\mathrm{hinge} + 0.3\,\mathrm{push} +
    0.2\,\mathrm{squat}$. These three are the only rows \texttt{GUIDE.md}
    gives explicitly; every other row is this project's draft, reviewed
    separately from the constants in this document (see
    \texttt{docs/AUDIT.md}'s closing note).
  \item $R_e : \mathcal{R} \to \mathbb{R}_{\ge 0}$, the \textbf{resource-cost
    row}, over
    $\mathcal{R} = \{\mathrm{knee}, \mathrm{post}, \mathrm{grip},
    \mathrm{shldr}, \mathrm{lats}, \mathrm{trunk},
    \mathrm{hang}\}$ (7 resources --- \texttt{cns} and \texttt{aero} both
    dropped, \texttt{hang} added; see Section~\ref{sec:score}'s reversal
    note). \texttt{aero} went in version 9: it had been authored for
    exactly four exercises, all monostructural with an empty $E$ row, so
    it was collinear with ``is a monostructural exercise'' --- a family
    label rather than a cost. Its removal left running and double-unders
    with empty $R$ rows, which the $\|R_e\| = 0$ guard in interference
    silently turns into zero similarity with everything, so real tissue
    costs were authored in its place. \texttt{hang} is capsule/tendon
    fatigue from hanging and swinging from a bar or rings (kipping
    pull-ups, toes-to-bar, muscle-ups, rope climbs), kept separate from
    \texttt{shldr} (pressing tension) and \texttt{grip} (crushing/holding
    load) since the injury mechanism differs from either. $R_e$ itself is
    \emph{not} normalized --- it is spend,
    not membership, so a heavier session costs more in absolute terms;
    only its use inside interference (below) is normalized.
  \item $A_e \subseteq \mathcal{S}$, the \textbf{admissible styles} --- which
    of light/heavy/hhr the exercise can be used for. This is what generates
    the cell mask below: a cell nothing can fill never accrues deficit.
  \item a \textbf{laterality} in $\{\mathrm{bilateral}, \mathrm{unilateral}\}$
    and an \textbf{implement} (barbell, dumbbell, kettlebell, bodyweight,
    \ldots), used for candidate selection and logging, not for scoring
    directly.
\end{itemize}

Neither $E$ nor $R$ is ever fit from data --- they are authored judgment,
versioned, held fixed. A stable-but-wrong map preserves relative trends;
a map that drifts does not.

\section{The log and exposures}

Every logged set carries a date, a style, an exercise id, and (for
unilateral exercises) a side. An \textbf{exposure event} is a distinct
$(\text{date}, s, e) \in \text{Date} \times \mathcal{S} \times \mathcal{E}$
triple --- multiple sets of the same exercise on the same day (a $5\times5$)
collapse to one event, and a unilateral exercise logged on both sides on
the same day also collapses to one event, not two.

\begin{fixedbox}
The implementation originally summed $E_{e,p}$ once per logged \emph{set
row} rather than once per exposure event, so a $5\times5$ counted as five
exposures instead of one, and a unilateral exercise logged L+R counted
double. Fixed by collapsing to the event definition above before applying
pattern weights. Two regression tests guard this
(\texttt{test\_multiple\_sets\_same\_day\_count\_as\_one\_exposure},
\texttt{test\_unilateral\_sides\_same\_day\_count\_as\_one\_exposure}).
\end{fixedbox}

Fix a trailing window of $W$ days ending on the query day $t$:
\[
X_{p,s}(t) \;=\! \sum_{\substack{(\text{date}, s', e)\ \text{exposure event} \\ s' = s,\ \text{date} \in (t-W,\, t]}} \!\! E_{e,p}
\]
the pattern-weighted count of exposures to pattern $p$ at style $s$ within
the window.

\begin{auditbox}
\texttt{WINDOW\_DAYS} $= 42$. \texttt{GUIDE.md} calls this ``the most
consequential single number'' in the whole model, but doesn't
second-guess the value itself. Every deficit and staleness figure below
is relative to this window, so it is worth confirming it still matches
your actual training rhythm before trusting anything downstream.
\end{auditbox}

\section{Deficits}

Each style has a weekly \textbf{target rate} $\tau_s$ (exposures per week
per pattern), each pattern a dimensionless \textbf{pattern weight}
$\omega_p$, and the target for cell $(p,s)$ over the window is
$\tau_s\, \omega_p\, w$ where $w = W/7$ is the window length in weeks.

\begin{center}
\begin{tabular}{lccl}
\toprule
$s$ & $\tau_s$ (per week per pattern) & source & note\\
\midrule
light & $1.35$ & \texttt{GUIDE.md}, trust & residual style, targeted highest\\
heavy & $0.50$ & \texttt{GUIDE.md}, trust & goal: heavy once per fortnight\\
hhr & $0.15$--$0.25$ & \texttt{GUIDE.md}, free choice & this code uses $0.20$\\
\bottomrule
\end{tabular}
\end{center}

\subsection{Pattern weights}
\label{sec:patternweight}

\[
\omega_p = \begin{cases}
0.25 & p = \mathrm{brace}\\
1 & \text{otherwise.}
\end{cases}
\]

Until dictionary version 8 targets were keyed by style alone and applied
uniformly across patterns, i.e.\ $\omega_p \equiv 1$. That was harmless
while $|\mathcal{P}|$ was fixed, and became load-bearing the moment it
wasn't: splitting \texttt{midline} into \texttt{flexion} and
\texttt{brace} under $\omega_p \equiv 1$ would have taken total weekly
demand from
\[
|\mathcal{P}| \sum_s \tau_s = 5 \times 2.05 = 10.25
\quad\text{to}\quad
6 \times 2.05 = 12.3
\]
against a supply of roughly $12$ pattern-bearing pieces per week --- the
entire capacity slack, spent silently by a taxonomy change that looks
purely descriptive. With $\omega_{\mathrm{brace}} = 0.25$, demand is
\[
\sum_{p}\omega_p \sum_s \tau_s = (5 + 0.25) \times 2.05 = 10.76,
\]
and the slack survives --- though only until version 9 added
$\mathrm{engine}$, which has no heavy cell but does have light and hhr
ones, taking demand to $12.31$ against the same $\approx 12$. The slack is
now slightly negative; see the audit box below.

\begin{auditbox}
$\omega_{\mathrm{engine}} = 1$ leaves total demand at $12.31$ per week
against a supply of roughly $12$ pattern-bearing pieces --- i.e.\ the
model is now \textbf{over-subscribed}, and the closed-loop simulation
shows it: \texttt{push/light} and \texttt{squat/light} finish the run
\emph{behind} target, the first positive end-state deficits any
simulation has produced. Lower $\omega_{\mathrm{engine}}$, raise the
supply estimate on the argument that a run does not compete for a lifting
slot, or accept over-subscription deliberately.
\end{auditbox}

\begin{auditbox}
$\omega_{\mathrm{brace}} = 0.25$ is a \textbf{free choice with nothing
behind it but the arithmetic above} --- it was chosen to fit the capacity
budget, not derived from a stated goal about how much carrying and
bracing is wanted. It is also, as of this writing, \emph{not what the
model actually delivers}: the 7-day any-exposure clock imposes a floor of
roughly one exposure per week on every pattern, which outvotes
$\omega_{\mathrm{brace}}$ entirely. See the audit box in
Section~\ref{sec:overdue}.
\end{auditbox}

\begin{auditbox}
\texttt{TARGET\_RATE["hhr"]} $= 0.20$. \texttt{GUIDE.md} gives a range
($0.15$--$0.25$, ``free choice; $0.25 \approx 2{:}1$ heavy:hhr'') and
leaves the exact value to you. This is the midpoint, not a derived
number --- pick your own ratio.
\end{auditbox}

A cell $(p,s)$ is \textbf{masked} (excluded from everything below) unless
some exercise can fill it:
\[
\text{admissible}(p,s) \iff \exists\, e \in \mathcal{E} : E_{e,p} > 0 \text{ and } s \in A_e.
\]

For admissible cells, the deficit is
\[
D_{p,s}(t) = \tau_s\, \omega_p\, w \;-\; \widetilde X_{p,s}(t), \qquad
\widetilde X_{p,s}(t) = \begin{cases}
X_{p,\mathrm{heavy}}(t) + \eta \cdot X_{p,\mathrm{hhr}}(t) & s = \mathrm{heavy}\\
X_{p,s}(t) & \text{otherwise,}
\end{cases}
\]
where $\eta$ is the \textbf{hhr credit}: hhr work partially, but never
fully, substitutes for fresh heavy work toward the heavy cell's own
deficit --- hhr's own cell is unaffected by $\eta$ and is scored against
its own target independently.

\begin{auditbox}
$\eta = \texttt{HHR\_CREDIT} = 0.5$. \texttt{GUIDE.md} gives both the
value and the physiological reasoning explicitly (loaded conditioning at
$60$--$70\%$ under fatigue trains work capacity and strength-endurance
but contributes little to maximal strength, which needs proximity to max
plus a rested nervous system), so this is probably right as specified.
Flagged anyway because \texttt{GUIDE.md} itself says ``if the priority
ever shifts toward maximal strength, lower this credit'' --- that call
belongs to you, not the code.
\end{auditbox}

Targets \emph{rank}; they are not reachable goals to report distance
from. With demand near capacity some cells always show shortfall by
design, and that is fine --- only the ordering needs to be right.

\section{Overdue overrides: the guarantee layer}
\label{sec:overdue}

Deficits rank; they do not enforce a hard floor on their own. A separate
guarantee layer gives a cell past a frequency clock a bonus large enough
to outrank everything else in the score:

\begin{center}
\begin{tabular}{llll}
\toprule
clock & threshold & satisfied by & applies to cell $(p,s)$ when\\
\midrule
any exposure & 7 days & any style & always\\
loaded & 14 days & heavy or hhr & $s \in \{\mathrm{heavy}, \mathrm{hhr}\}$, $p \notin \mathcal{X}$\\
classic heavy & 24 days & heavy only & $s = \mathrm{heavy}$, $p \notin \mathcal{X}$\\
\bottomrule
\end{tabular}
\end{center}

Formally, with $\mathrm{last}(p, \mathcal{S}') $ the most recent exposure
date to pattern $p$ at any style in $\mathcal{S}' \subseteq \mathcal{S}$
(or $-\infty$ if never):
\[
\mathrm{overdue}(p,s,t) = \begin{cases}
B & t - \mathrm{last}(p,\mathcal{S}) > 7\\
B & p \notin \mathcal{X},\ s \in \{\mathrm{heavy},\mathrm{hhr}\},\ t - \mathrm{last}(p,\{\mathrm{heavy},\mathrm{hhr}\}) > 14\\
B & p \notin \mathcal{X},\ s = \mathrm{heavy},\ t - \mathrm{last}(p,\{\mathrm{heavy}\}) > 24\\
0 & \text{otherwise,}
\end{cases}
\]
checked in that order (the first satisfied condition wins), where
$\mathcal{X} = \{\mathrm{flexion}, \mathrm{brace}, \mathrm{engine}\}$ is
the set of \textbf{clock-exempt} patterns. Both are exempt from the 14- and 24-day
clocks --- every other pattern loads the trunk incidentally, so a
dedicated heavy-trunk target would overstate the need, and \texttt{brace}
additionally consists of movements that are loaded by nature and would
otherwise trip the classic-heavy clock unprompted --- but \emph{not} from
the 7-day any-exposure clock.

\begin{auditbox}
That last exemption boundary is currently doing something unintended.
The 7-day clock implies a floor of roughly one exposure per week, while
$\omega_{\mathrm{brace}} = 0.25$ asks for about half that. The clock
always wins, so \texttt{brace}'s effective rate is about twice its stated
target. Either $\mathcal{X}$'s members need a longer any-exposure
threshold, or $\omega_{\mathrm{brace}}$ should rise to match a weekly
floor. See \texttt{docs/AUDIT.md}.
\end{auditbox}
The loaded and classic-heavy clocks only bonus the styles that can
actually resolve them: a light pick can't clear a ``loaded'' debt, so
$\mathrm{overdue}(p, \mathrm{light}, t)$ never receives the 14- or 24-day
bonus, only the 7-day one.

\begin{auditbox}
$B = \texttt{OVERDUE\_BONUS} = 1000.0$. \texttt{GUIDE.md} specifies only
that the bonus must be ``large enough to outrank everything''; this is a
sentinel constant, not a bound derived from the other terms' ranges. Safe
as long as no plausible combination of the terms in Section~\ref{sec:score}
can approach it, but worth re-checking if the $\lambda$'s or staleness
caps below are ever changed by a large factor.
\end{auditbox}

\section{Candidate selection}
\label{sec:selection}

For an admissible cell $(p,s)$, the eligible exercises are
$\{e \in \mathcal{E} : E_{e,p} > 0,\ s \in A_e\}$. The candidate offered
for that cell is the \textbf{stalest eligible variant}:
\[
\hat e(p,s,t) = \operatorname*{arg\,max}_{e \text{ eligible}} \Big(t - \mathrm{lastlogged}(e),\ E_{e,p},\ e\Big)
\]
lexicographically --- staleness first (never-logged treated as
$+\infty$), pattern weight as tie-break, exercise id as final tie-break
for determinism. Because each cell contributes only its stalest variant,
the second-best option for that cell is never shown --- a known,
documented limitation, not a bug.

\section{Scoring}
\label{sec:score}

For each admissible cell, with $\hat e = \hat e(p,s,t)$ as the candidate
and $C$ the set of exercises already chosen this session:
\begin{align}
\mathrm{score}(p,s) \;=\;\; & \widehat D_{p,s}(t) \notag\\
+\;\; & \gamma_p \cdot M_p(t) \notag\\
+\;\; & \lambda_{\mathrm{cell}} \cdot \frac{\min\!\big(t - \mathrm{last}(p,\{s\}),\ 21\big)}{21} \notag\\
+\;\; & \lambda_{\mathrm{uni}} \cdot \frac{\min\!\big(\max(u_p, 0),\ 3.0\big)}{3.0} \notag\\
+\;\; & \mathrm{overdue}(p,s,t) \notag\\
-\;\; & 0.25 \sum_{c \in C} \widehat R_{\hat e} \, W \, \widehat R_c^{\!\top} \notag\\
-\;\; & 1.1 \sum_{c \in C} \mathbb{1}\!\left[\exists\, p' : E_{\hat e,p'} > 0 \wedge E_{c,p'} > 0\right] \notag
\end{align}
where $\widehat R_e = R_e / \|R_e\|_2$ is $e$'s resource-cost row unit-normalized.
with the following pieces.

\paragraph{Reversal: normalized deficit and staleness, $\lambda_{\mathrm{move}}$ removed.}
$\widehat D_{p,s}(t) = D_{p,s}(t) / (\tau_s \omega_p w)$ --- the deficit
expressed as a fraction of that \emph{cell's} own target. Since exposure
is never negative, $D_{p,s} \le \tau_s \omega_p w$ always, so $\widehat D$
is bounded above by $1$ with no new machinery. Dividing by the
pattern-scaled target (rather than $\tau_s w$) extends the same guarantee
across patterns now that $\omega_p \not\equiv 1$: a fully-neglected
\texttt{brace} cell saturates at $1$ exactly like a fully-neglected
\texttt{squat} cell, despite an absolute target a quarter the size.
Without it \texttt{brace} could never outrank \texttt{squat} no matter how
long it went untouched --- the identical failure this normalization was
introduced to fix across styles. This matters because targets vary a lot in
absolute size (hhr's is $1.2$ over the window, light's is $8.1$):
comparing raw deficits structurally favours whichever style has the
biggest target, exactly the problem \texttt{recommend()} (Section~\ref{sec:recommend})
already had to correct for one level up. Normalizing fixes it here too: a
fully-neglected hhr cell and a fully-neglected light cell both saturate
at $\widehat D = 1$, on equal footing. Cell staleness and the unilateral
shortfall are divided by their own caps ($21$ days, $3.0$) the same way,
so every additive term above lives on a comparable $[0,1]$-ish scale, and
the $\lambda$'s are weights on that scale rather than magnitude-matched
placeholders guessed against a moving target.

$\lambda_{\mathrm{move}}$ (movement-level staleness) has been removed
entirely. It never affected \emph{which} exercise got offered for a cell
--- that was always $\hat e(p,s,t)$ from Section~\ref{sec:selection},
independent of scoring. Its only job was a secondary nudge to a cell's
ranking when its offered candidate happened to be stale even though the
cell overall wasn't; GUIDE.md records that an ablation on the discarded
prototype had already found it moved outcomes by $\sim 0.05$ on one seed
--- removing it is a real simplification, not a loss of a load-bearing
mechanism.

\emph{Deliberately not done:} no unbounded/accumulating staleness term
was added to compensate. An accumulating deficit (never capped, only
reset by servicing, à la Deficit Round Robin scheduling) would give the
same eventual-guarantee property the overdue clocks below provide, but
this model does greedy top-1 selection rather than round-robin service,
so an unbounded backlog could monopolize the schedule for an extended
stretch once large enough. The clocks stay a separate, bounded, fixed-
magnitude override rather than a growing one.

\begin{auditbox}
\textbf{Correction from an earlier version of this document:}
$\gamma_{\mathrm{behind}}=0.35$/$\gamma_{\mathrm{over}}=0.10$ and the
resource-overlap coefficient $0.25$ below were previously described as
trusted GUIDE.md values, reasoning that GUIDE.md states their numbers
directly. That conflated ``has a specific number attached'' with
``derived from stated goals.'' GUIDE.md's own parameter classification
places \emph{both} in the tuned-against-synthetic-history, near-noise
class --- same as the $\lambda$'s. The asymmetry in $\gamma$ (behind
pulls harder than over-covered pushes back) is the load-bearing idea;
the specific magnitudes came from the same ablation process as everything
else in this box, not from GUIDE.md's stated requirements.

$\lambda_{\mathrm{cell}} = 0.25$ (was $0.10$), $\lambda_{\mathrm{uni}} =
0.10$ (was $0.30$) --- rescaled for the normalization above, by
preserving each term's old ratio to deficit's old maximum: $\lambda_{\mathrm{cell}}$
contributed up to $0.10 \times 21 = 2.1$ against a max deficit of $8.1$,
a ratio of $\approx 0.26$; $\lambda_{\mathrm{uni}}$ contributed up to
$0.30 \times 3.0 = 0.9$ against the same $8.1$, a ratio of $\approx
0.11$. Rounded to clean numbers rather than false precision. Still a
guess, just a rescaled one --- GUIDE.md's own posture here is ``leave
conservative until real logs exist.''

\textbf{Newly discovered, not yet resolved:} the resource-overlap
coefficient ($0.25$, below) and the same-pattern-repeat coefficient
($1.1$) were left at their original values, since neither was part of
what this normalization pass touched. But both multiply a \emph{raw},
unnormalized magnitude ($R\,W\,R^\top$, and a $0/1$ indicator), the same
kind of quantity $\lambda_{\mathrm{cell}}$ used to be calibrated against
before it was rescaled. With deficit now capped at $1$ instead of running
up to $8$, these two penalty terms are relatively far more powerful than
before, and can now push a candidate's score negative on their own (the
worked example in Section~\ref{sec:example} hits exactly this). Whether
that's a problem or a reasonable consequence of "resource conflicts
should genuinely veto a pick" is an open question, not yet decided.
\end{auditbox}

\paragraph{Coupling across styles.} $M_p(t)$ is now the \emph{mean} of
$\widehat D_{p,s'}(t)$ over $p$'s admissible styles $s'$ (not the sum, as
it was before normalization) --- keeping it on the same $[-\infty,1]$-ish
scale as $\widehat D$ itself:
\[
M_p(t) = \frac{1}{|\{s' : \text{admissible}(p,s')\}|} \sum_{s' \in \mathcal{S}} \widehat D_{p,s'}(t),
\qquad
\gamma_p = \begin{cases} 0.35 & M_p(t) > 0\\ 0.10 & \text{otherwise.}\end{cases}
\]
The asymmetry matters: without it, over-coverage in one style starved
patterns that had few admissible cells in another.

\paragraph{Unilateral shortfall.} $u_p = \texttt{UNILATERAL\_TARGET}
\cdot w - U_p(t)$, where $\texttt{UNILATERAL\_TARGET} = 0.30$ per week per
pattern (\texttt{GUIDE.md}, trust) and $U_p(t)$ is the windowed,
pattern-weighted count of exposure events from unilateral exercises,
clamped to $[0, 3.0]$ before normalizing --- uncapped, every pattern went
maximally behind at once the first time this was implemented, and the
whole candidate list went unilateral.

\paragraph{Reversal: resource interference is now cosine similarity.}
$W \in \mathbb{R}^{8\times8}$ is the resource-interaction matrix, so
$\widehat R_{\hat e}\,W\,\widehat R_c^{\top}$ --- with $W = I$, this
reduces to the ordinary cosine similarity between $\hat e$ and $c$'s
resource-cost profiles, bounded to $[0,1]$ since no resource cost is
negative. This is what lets the model know a bike ride blocks a heavy
squat despite sharing no movement pattern.

Two changes from an earlier version of this document, made together.
First, \texttt{cns} was dropped from $\mathcal{R}$ entirely: it was
present in every one of the 64 exercises, making
it the single most universally-shared resource and the dominant driver of
interference between almost any two picks, regardless of whether they
actually competed for the same tissue --- also the least concretely
groundable of the eight numbers being authored per exercise to begin
with. Second, interference used to be the \emph{raw} dot product
$R_{\hat e}\,W\,R_c^\top$, unbounded and scaling with both vectors'
absolute magnitude; once deficit was normalized to $[0,1]$, this became
disproportionately powerful, as Section~\ref{sec:example} showed directly
(a candidate's score flipped negative purely from resource overlap).
Unit-normalizing $R$ before the dot product --- i.e. cosine similarity
--- fixes that by construction, and also changes what the term
\emph{means}: from ``how much absolute stress do these share'' to ``how
similar are their stress profiles,'' which is arguably the more correct
notion of interference independent of the rescaling motivation.

\begin{auditbox}
\texttt{GUIDE.md} defines interference as $R\,W\,R^{\top}$ but never gives
$W$'s values. It currently defaults to the $8\times 8$ identity matrix ---
no assumed cross-resource coupling; \texttt{grip} fatigue doesn't spill
into \texttt{shldr} fatigue, \texttt{hang} doesn't spill into either, and
so on. That is the most conservative possible choice, not a reasoned one.
\texttt{hang} is the resource most likely to actually deserve coupling to
\texttt{shldr}/\texttt{grip} (a kipping shoulder is a real overuse
pattern, not an independent one) --- worth a look once real logs exist.
If some resources should interact,
edit the matrix directly in \texttt{exercises.json}. Separately,
\texttt{RESOURCE\_OVERLAP\_COEF} $= 0.25$ is unchanged in value but not in
meaning: it now weights a bounded $[0,1]$ quantity, so it reads as ``a
fully-overlapping resource profile costs as much as maximal cell
staleness'' ($=\lambda_{\mathrm{cell}}$) rather than a fraction of an
unbounded raw magnitude. Still a guess pending real logs.
\end{auditbox}

\paragraph{Same-pattern repeat.} The indicator sum counts how many already-chosen
exercises this session share at least one pattern with the candidate.
This is a soft penalty, not a mask --- back squat then front squat stays
reachable if below-parallel is genuinely starved --- with coefficient
$1.1$, of unclear provenance: GUIDE.md's near-noise list doesn't name it
explicitly, but it isn't in the trust-these requirements table either.

\section{Hard and soft session rules}

One rule is a hard mask rather than a scoring penalty: a pattern loaded
today (by any already-chosen exercise, at any style) cannot reappear at
\texttt{hhr} in the same session --- no heavy squat followed by squatting
at high heart rate. Formally, when generating candidates for
$s = \mathrm{hhr}$, the cell $(p, \mathrm{hhr})$ is dropped entirely
(not merely penalized) whenever $p \in \bigcup_{c \in C}\{p' : E_{c,p'}>0\}$.

Every other interaction between already-chosen exercises and a new
candidate --- same pattern at heavy, shared resource cost --- is a soft
penalty already folded into the score in Section~\ref{sec:score}, not a
second masking layer.

\section{Rep scheme and rep count}

Rep schemes apply to \texttt{heavy} only (an \texttt{hhr} piece gets a rep
target, not a scheme). For a pattern $p$ due for heavy work, the
suggested scheme is whichever entry in a fixed menu was least recently
used for that pattern at \texttt{heavy} --- ranked per pattern, not
globally, so heavy cleans avoid schemes recently used by heavy deadlifts.

\begin{auditbox}
The scheme menu (\texttt{REP\_SCHEMES}: \texttt{5x5}, \texttt{3x3x3x3x3},
\texttt{5-3-1}, \texttt{8x3}, \texttt{3RM build}, \texttt{10x2},
\texttt{4x6}) doesn't come from \texttt{GUIDE.md} --- it describes
schemes being ranked per pattern by staleness but never names an actual
menu. This list is a generic placeholder; swap in whatever you actually
program.
\end{auditbox}

Rep counts are not prescribed at all: \texttt{suggest\_rep\_count} returns
the rolling median of reps logged for a given (exercise, style) pair, or
\texttt{None} if there is no history yet. Rep counts are learned from the
log, never authored, because rep counts are skewed and a median is the
right summary; schemes are categorical and cannot be averaged the same
way, which is why they are ranked instead.

\section{The \texttt{recommend()} verdict}
\label{sec:recommend}

\texttt{recommend()} answers which \emph{mode} (not cell) is most needed,
comparing styles \textbf{relative to each style's own target} --- light is
targeted at nearly three times heavy's rate, so comparing raw deficits
would always favor light regardless of actual need. For each style with
at least one admissible cell:
\[
\rho_s(t) = 100 \cdot \frac{\sum_{p} D_{p,s}(t)}{\sum_{p : \text{admissible}(p,s)} \tau_s\, \omega_p\, w}
\]
the aggregate deficit as a percentage of that style's own aggregate
target. Let $s^\star = \operatorname{arg\,max}_s \rho_s(t)$.

\begin{fixedbox}
Ranking $s^\star$ purely by $\rho_{s^\star}(t) = \max_s \rho_s(t)$ without
checking its sign let the verdict declare a style ``needed'' when
\emph{every} style was already over its target. A two-month closed-loop
simulation showed exactly this: once every $\rho_s(t)$ went negative
(all styles over-covered — an artifact of the simulation forcing a fixed
number of pieces per session with no stopping rule), light kept "winning"
the comparison every week, with the reported margin \emph{growing},
purely because heavy and hhr were becoming more over-covered
\emph{faster} --- not because light was falling behind. Fixed: if
$\rho_{s^\star}(t) \le 0$, the verdict is unconditionally
\texttt{no strong signal} with no style named, regardless of the margin
to the runner-up. Guarded by
\texttt{test\_recommend\_does\_not\_favour\_the\_least\_overshot\_style}.
\end{fixedbox}

When $\rho_{s^\star}(t) > 0$, let $\rho_{(2)}(t)$ be the second-highest
value and $\mu(t) = \rho_{s^\star}(t) - \rho_{(2)}(t)$ the margin. The
verdict bands are
\[
\mathrm{band}(t) = \begin{cases}
\texttt{clear} & \mu(t) > 25\\
\texttt{leaning} & \mu(t) > 10\\
\texttt{no strong signal} & \text{otherwise,}
\end{cases}
\]
with the style named only in the first two bands --- \texttt{GUIDE.md} is
explicit that most days genuinely have no right answer, and the output
should say so rather than force a pick. All three band thresholds
($25$, $10$) are given directly by \texttt{GUIDE.md}, not tuned here.

\section{A worked example}
\label{sec:example}

Everything above is abstract; this section runs the actual pipeline on a
small, hand-picked slice of the dictionary and a short synthetic log, so
every number below can be checked by hand. This is illustrative only ---
a reduced 5-exercise dictionary, not the full dictionary itself, chosen
to touch every mechanism in the model with round numbers.

\begin{auditbox}
\textbf{This section is still on the pre-version-8 five-pattern set and
has not been reworked.} Everything above it reflects dictionary version 8
(six patterns, $\omega_p$, $\mathcal{X}$); the arithmetic below was
hand-checked against version 7 and is left intact rather than
half-migrated, because the two are not reconcilable by renaming: the
example's fifth exercise is a kettlebell swing carrying
$\mathrm{midline} = 0.2$, and version 8 removed exactly that credit on the
grounds that a swing's reactive bracing is a resource cost, not
adaptation-specific transfer. Reworking it means choosing a new fifth
exercise and recomputing the deficit table, the normalized deficits and
both scored cells --- worth doing deliberately, not as a search-replace.
Read this section as a correct illustration of the \emph{mechanisms},
with stale pattern names.
\end{auditbox}

\subsection{A reduced dictionary}

Five exercises, one per pattern (kettlebell swing also touches midline as
a secondary pattern, showing what a nonzero residual looks like):

\begin{center}
\begin{tabular}{lccccccl}
\toprule
exercise & hinge & push & squat & pull & midline & residual & admissible styles\\
\midrule
back squat & -- & -- & 1.0 & -- & -- & 0.0 & heavy, hhr\\
kettlebell swing & 0.4 & -- & -- & -- & 0.2 & 0.4 & light, hhr\\
strict pull-up & -- & -- & -- & 1.0 & -- & 0.0 & light, heavy, hhr\\
strict press & -- & 1.0 & -- & -- & -- & 0.0 & heavy, hhr\\
deadlift & 1.0 & -- & -- & -- & -- & 0.0 & light, heavy, hhr\\
\bottomrule
\end{tabular}
\end{center}
\begin{center}\small this is the matrix $E$ --- rows are exercises, columns are patterns, ``--'' is 0\end{center}

\begin{center}
\begin{tabular}{lccccccc}
\toprule
exercise & knee & post & grip & shldr & lats & trunk & hang\\
\midrule
back squat & 3 & 1 & -- & -- & -- & 2 & --\\
kettlebell swing & -- & 1 & 1 & -- & -- & 1 & --\\
strict pull-up & -- & -- & 1 & -- & 2 & -- & 1\\
strict press & -- & -- & -- & 3 & -- & 1 & --\\
deadlift & -- & 3 & 2 & -- & -- & 2 & --\\
\bottomrule
\end{tabular}
\end{center}
\begin{center}\small this is the matrix $R$ (\texttt{aero} omitted --- none of these five touch it; \texttt{cns} is gone from the model entirely, see Section~\ref{sec:score}. Strict pull-up's \texttt{hang}$=1$ doesn't feed into any calculation shown below --- none of the interference examples involve it.)\end{center}

All five are bilateral, so the unilateral-shortfall term below comes out
identical across patterns in this example --- dictionary version 7 has
unilateral variants for all of these; none are included here to keep the
arithmetic small. $W$, the resource-interaction matrix, is the
$8\times8$ identity throughout (Section~\ref{sec:score}'s audit box) ---
with $W=I$, $\widehat R_a\,W\,\widehat R_b^\top$ reduces to ordinary
cosine similarity, so it is not worth drawing out as a $64$-entry table
of mostly zeros.

From this dictionary alone, three cells are masked before any log exists
at all: \textbf{(squat, light)} --- no light-admissible squat exercise
here --- and likewise \textbf{(push, light)} and \textbf{(midline, heavy)}.
(Dictionary version 7 does have light-squat and heavy-midline options; the
mask here is an artifact of using only 5 exercises, not a property of the
real dictionary.) 12 cells remain admissible.

\subsection{A five-event log and the resulting deficits}

Take the query day $t$ as ``today,'' with the usual $42$-day window
($w=6$ weeks). The log:

\begin{center}
\begin{tabular}{llll}
\toprule
days before $t$ & exercise & style & contributes\\
\midrule
3 & back squat & heavy & $\mathrm{squat} \mathrel{+}= 1.0$\\
5 & kettlebell swing & hhr & $\mathrm{hinge} \mathrel{+}= 0.4,\ \mathrm{midline} \mathrel{+}= 0.2$\\
15 & strict pull-up & heavy & $\mathrm{pull} \mathrel{+}= 1.0$\\
20 & deadlift & light & $\mathrm{hinge} \mathrel{+}= 1.0$\\
25 & strict press & heavy & $\mathrm{push} \mathrel{+}= 1.0$\\
\bottomrule
\end{tabular}
\end{center}

All five fall inside the window (the oldest, at day 25, is well clear of
the day-42 edge, picked deliberately so none of the near-the-boundary
cliff behaviour discussed earlier shows up here and complicates the
arithmetic). Windowed exposures $X_{p,s}$, everywhere else $0$:
\begin{align*}
X_{\mathrm{squat,heavy}} &= 1.0, & X_{\mathrm{hinge,hhr}} &= 0.4, & X_{\mathrm{midline,hhr}} &= 0.2,\\
X_{\mathrm{pull,heavy}} &= 1.0, & X_{\mathrm{hinge,light}} &= 1.0, & X_{\mathrm{push,heavy}} &= 1.0.
\end{align*}

Targets: $\tau_{\mathrm{light}}\,w = 1.35 \times 6 = 8.1$,
$\tau_{\mathrm{heavy}}\,w = 0.50 \times 6 = 3.0$,
$\tau_{\mathrm{hhr}}\,w = 0.20 \times 6 = 1.2$. Only hinge's heavy cell picks
up the $\eta=0.5$ hhr credit, from the kettlebell-swing exposure:
$\widetilde X_{\mathrm{hinge,heavy}} = 0 + 0.5 \times 0.4 = 0.2$; every other
heavy cell here has zero hhr exposure to credit, so $\widetilde X = X$
for them.

\begin{center}
\begin{tabular}{lrrr}
\toprule
pattern & light & heavy & hhr\\
\midrule
hinge & $8.1 - 1.0 = 7.1$ & $3.0 - 0.2 = 2.8$ & $1.2 - 0.4 = 0.8$\\
push & \textit{masked} & $3.0 - 1.0 = 2.0$ & $1.2 - 0 = 1.2$\\
squat & \textit{masked} & $3.0 - 1.0 = 2.0$ & $1.2 - 0 = 1.2$\\
pull & $8.1 - 0 = 8.1$ & $3.0 - 1.0 = 2.0$ & $1.2 - 0 = 1.2$\\
midline & $8.1 - 0 = 8.1$ & \textit{masked} & $1.2 - 0.2 = 1.0$\\
\bottomrule
\end{tabular}
\end{center}
\begin{center}\small $D_{p,s}(t)$, the full deficit table for this state\end{center}

\subsection{Normalized deficits and the coupling term}

Dividing the table above by each column's own target ($8.1$/$3.0$/$1.2$)
gives $\widehat D_{p,s}(t)$, and averaging each pattern's admissible
entries gives $M_p(t)$ (Section~\ref{sec:score}):

\begin{center}
\begin{tabular}{lrrr r}
\toprule
pattern & light & heavy & hhr & $M_p(t)$\\
\midrule
hinge & $0.877$ & $0.933$ & $0.667$ & $0.826$\\
push & \textit{masked} & $0.667$ & $1.000$ & $0.833$\\
squat & \textit{masked} & $0.667$ & $1.000$ & $0.833$\\
pull & $1.000$ & $0.667$ & $1.000$ & $0.889$\\
midline & $1.000$ & \textit{masked} & $0.833$ & $0.917$\\
\bottomrule
\end{tabular}
\end{center}

Every $M_p(t) > 0$ here, so $\gamma_p = 0.35$ throughout this example.

\subsection{Scoring two cells: overdue versus not}

\textbf{Cell (push, heavy).} The only pattern-push exercise here, strict
press, is also the only eligible candidate, so $\hat e = \text{strict
press}$ regardless of staleness. Pattern push was last touched (any style)
25 days ago, so $t - \mathrm{last}(\mathrm{push},\mathcal{S}) = 25 > 7$: the
7-day clock alone is already tripped, so $\mathrm{overdue} = 1000$ before
the 14- and 24-day clocks are even consulted. With nothing chosen yet
this session ($C = \emptyset$, so the overlap and repeat terms are both
$0$):
\begin{align*}
\mathrm{score}(\mathrm{push},\mathrm{heavy})
&= \underbrace{0.667}_{\widehat D} + \underbrace{0.35 \times 0.833}_{\gamma M = 0.292}
 + \underbrace{0.25 \times \tfrac{\min(25,21)}{21}}_{\lambda_{\mathrm{cell}} = 0.250}\\
&\quad + \underbrace{0.10 \times \tfrac{1.8}{3.0}}_{\lambda_{\mathrm{uni}} = 0.060}
 + \underbrace{1000}_{\mathrm{overdue}}\\
&= 1001.27.
\end{align*}
(The unilateral term is $0.10 \times \min(\max(0.30\times 6 - 0, 0),
3.0)/3.0 = 0.10 \times 0.6 = 0.060$ for every pattern here, since none of
the five exercises are unilateral.)

\textbf{Cell (hinge, light).} Two exercises are eligible --- deadlift and
kettlebell swing both allow light and both carry hinge weight. Deadlift was
last logged 20 days ago; kettlebell swing 5 days ago, \emph{as hhr} ---
which no longer matters for scoring since $\lambda_{\mathrm{move}}$ is
gone, but it still decides which exercise selection offers, and deadlift
is stalest. Pattern hinge was touched 5 days ago (that same
kettlebell-swing set), so $t - \mathrm{last}(\mathrm{hinge},\mathcal{S}) = 5
\le 7$: not overdue --- and since this cell's style is \texttt{light},
the 14- and 24-day clocks never apply to it regardless of staleness
(Section~\ref{sec:overdue}):
\begin{align*}
\mathrm{score}(\mathrm{hinge},\mathrm{light})
&= \underbrace{0.877}_{\widehat D} + \underbrace{0.35 \times 0.826}_{0.289}
 + \underbrace{0.25 \times \tfrac{\min(20,21)}{21}}_{0.238}\\
&\quad + \underbrace{0.10 \times \tfrac{1.8}{3.0}}_{0.060}
 + \underbrace{0}_{\mathrm{overdue}}\\
&= 1.46.
\end{align*}

Even though (hinge, light) has the second-largest raw deficit in the whole
table, (push, heavy) outranks it by roughly $685\times$ --- \emph{more}
lopsided than before normalization (it was $\sim70\times$), because
normalizing shrank the non-overdue scores without touching the overdue
bonus at all. Exactly what ``a bonus large enough to outrank everything''
is supposed to do, and a reminder that $B=1000$ has even more headroom
now (see Section~\ref{sec:overdue}'s audit box). (Not shown in detail:
(hinge, heavy) is \emph{also} overdue in this state, but only via the
24-day classic-heavy clock --- hinge has never actually been trained at
true heavy style in this log, so
$t - \mathrm{last}(\mathrm{hinge},\{\mathrm{heavy}\}) = \infty$, even
though the recent hhr exposure keeps both the 7- and 14-day clocks
satisfied. This is the reason the model needs three separate clocks
rather than one: hhr work can keep a pattern ``warm'' without ever
satisfying the specific requirement for fresh heavy work.)

\subsection{A second pick: resource overlap, now properly scaled}

Suppose (push, heavy) is chosen first, so $C = \{\text{strict press}\}$.
Now score (squat, hhr) --- candidate: back squat, the only squat exercise
admissible at hhr. With $\mathrm{cns}$ gone, the two resource-cost rows
are $R_{\text{back squat}} = (\mathrm{knee}{:}3, \mathrm{post}{:}1,
\mathrm{trunk}{:}2)$ and $R_{\text{strict press}} = (\mathrm{shldr}{:}3,
\mathrm{trunk}{:}1)$, sharing only $\mathrm{trunk}$:
\[
\widehat R_{\text{back squat}} \cdot W \cdot \widehat R_{\text{strict press}}^\top
= \frac{2 \times 1}{\|(3,1,2)\|\;\|(3,1)\|}
= \frac{2}{3.742 \times 3.162} = 0.169.
\]
Back squat and strict press share no pattern ($\mathrm{squat} \ne
\mathrm{push}$), so the same-pattern-repeat term is $0$. Pattern squat was
touched 3 days ago by that same back-squat set, which also satisfies
squat's own loaded clock, so (squat, hhr) is not overdue. $\widehat
D_{\mathrm{squat,hhr}} = 1.2/1.2 = 1.0$ (squat has never had an hhr exposure,
so this cell is fully unmet):
\begin{align*}
\mathrm{score}(\mathrm{squat},\mathrm{hhr})
&= \underbrace{1.0}_{\widehat D} + \underbrace{0.35 \times 0.833}_{0.292}
 + \underbrace{0.25 \times \tfrac{21}{21}}_{0.250} + \underbrace{0.10 \times \tfrac{1.8}{3.0}}_{0.060}\\
&\quad + 0
 - \underbrace{0.25 \times 0.169}_{0.042}
 - 0\\
&= 1.559.
\end{align*}
Comfortably positive, and this is the direct payoff of the two changes in
Section~\ref{sec:score}'s reversal: under the raw, pre-\texttt{cns}-removal
formula this same candidate scored $-0.398$ --- a $2$-point overlap
penalty (driven mostly by $\mathrm{cns}{:}3\times2=6$ of the raw dot
product's $8$) exceeded every positive term combined. With \texttt{cns}
gone and the remaining overlap unit-normalized, the penalty drops to
$0.042$: back squat and strict press turn out to share very little
resource cost once the artificially-universal \texttt{cns} entry is
removed, and the residual overlap costs a proportionate amount rather
than dominating the score.

This is also where the hard rule bites: (push, hhr), with a normalized
deficit of $0.8/1.2 = 0.667$, would ordinarily be a live candidate in the
workout pool --- but because push was already loaded today (by strict
press, at heavy), that cell is \emph{dropped from the candidate list
entirely}, not merely penalized. It cannot appear no matter how far
behind it is, for the rest of this session.

\subsection{The \texttt{recommend()} verdict on this state}

Using the same deficit table, at the original state (before either pick
above):
\begin{align*}
\rho_{\mathrm{light}} &= 100 \cdot \frac{7.1 + 8.1 + 8.1}{8.1 \times 3} = 95.9\%,\\
\rho_{\mathrm{heavy}} &= 100 \cdot \frac{2.8+2.0+2.0+2.0}{3.0 \times 4} = 73.3\%,\\
\rho_{\mathrm{hhr}} &= 100 \cdot \frac{0.8+1.2+1.2+1.2+1.0}{1.2 \times 5} = 90.0\%.
\end{align*}
All three are positive (nothing is over-covered), so the sign guard from
Section~\ref{sec:recommend} does not apply. Light ranks first, hhr
second, margin $95.9 - 90.0 = 5.9$ --- under the leaning threshold of
$10$, so the verdict is \texttt{no strong signal} with no style named,
even though light's raw deficit sum ($23.3$) is nearly double heavy's
($8.8$). That gap is exactly why the comparison is relative to each
style's own target rather than on raw deficits: light's target is also
nearly triple heavy's, and once that is accounted for, the three styles
are much closer to equally served than the raw sums suggest.

\section*{Summary: every item flagged above}

\begin{center}
\small
\begin{tabular}{p{3.8cm} p{2.2cm} p{8.5cm}}
\toprule
constant & value used & status\\
\midrule
\texttt{WINDOW\_DAYS} & 42 & GUIDE.md value, not second-guessed\\
\texttt{TARGET\_RATE["hhr"]} & 0.20 & midpoint of GUIDE.md's 0.15--0.25 range\\
\texttt{HHR\_CREDIT} & 0.5 & GUIDE.md value; ``lower if priority shifts'' is your call\\
\texttt{OVERDUE\_DAYS} & 7 / 14 / 24 & GUIDE.md values, stated as goals-as-constraints\\
\texttt{GAMMA\_BEHIND / GAMMA\_OVER} & 0.35 / 0.10 & tuned/near-noise per GUIDE.md's own classification (corrected --- previously mislabeled ``trust''); asymmetry is the point, magnitudes aren't\\
\texttt{LAMBDA\_CELL} & 0.25 (was 0.10) & rescaled for normalized [0,1] deficit; still a placeholder\\
\texttt{LAMBDA\_MOVE} & \emph{removed} & simplification; selection is unaffected, ablation found it near-noise\\
\texttt{LAMBDA\_UNI} & 0.10 (was 0.30) & rescaled for normalized [0,1] deficit; still a placeholder\\
\texttt{UNILATERAL\_\allowbreak SHORTFALL\_CAP} & 3.0 & unspecified in GUIDE.md; placeholder (defines what gets normalized, unchanged)\\
resource matrix $W$ & identity ($8\times8$) & unspecified in GUIDE.md; conservative default; \texttt{cns} row/col removed, \texttt{hang} row/col added\\
\texttt{RESOURCE\_OVERLAP\_COEF} & 0.25 & tuned/near-noise per GUIDE.md; same value, new meaning --- weights bounded cosine similarity (Section~\ref{sec:score}'s reversal), not raw magnitude\\
\texttt{SAME\_PATTERN\_\allowbreak REPEAT\_COEF} & 1.1 & provenance unclear --- not in GUIDE.md's near-noise list nor its trust-these table; still a raw count, not normalized by session size\\
\texttt{OVERDUE\_BONUS} & 1000.0 & sentinel, not derived from other terms; even more headroom now that other terms are normalized\\
\texttt{REP\_SCHEMES} & generic list & not specified in GUIDE.md; placeholder menu\\
band thresholds & 25 / 10 & GUIDE.md values\\
$E$/$R$ dictionary rows & \multicolumn{2}{p{10.5cm}}{authored draft beyond the 3 GUIDE.md examples --- separate review}\\
\bottomrule
\end{tabular}
\end{center}

\begin{sloppypar}
None of the placeholder rows above have been exercised against real
training data. Per \texttt{docs/\allowbreak GUIDE.md}'s own build
order, that is the next step, not more constant-tweaking in the abstract.
The resource-overlap coefficient's disproportionate power (an earlier
version of this document showed it flipping a candidate's score negative)
has since been resolved by dropping \texttt{cns} and normalizing
interference to cosine similarity, both described in Section~\ref{sec:score}
and demonstrated in Section~\ref{sec:example}. The same-pattern-repeat
coefficient's analogous issue --- a raw count rather than a normalized
fraction --- has not.
\end{sloppypar}

\end{document}
